% Generated draft: requires mathematical and source review.
\documentclass[11pt]{article}
\usepackage[margin=25mm]{geometry}
\usepackage{fontspec}
\setmainfont{Latin Modern Roman}
\usepackage{amsmath,amssymb,mathtools,unicode-math}
\setmathfont{Latin Modern Math}
\usepackage{enumitem,needspace,xurl,hyperref}
\setlength{\parindent}{0pt}
\setlength{\parskip}{5pt}
\setlength{\emergencystretch}{3em}
\newcommand{\N}{\mathbb{N}}
\newcommand{\Z}{\mathbb{Z}}
\newcommand{\Q}{\mathbb{Q}}
\newcommand{\R}{\mathbb{R}}
\newcommand{\C}{\mathbb{C}}
\newcommand{\F}{\mathbb{F}}
\newcommand{\A}{\mathbb{A}}
\newcommand{\PP}{\mathbb{P}}
\newcommand{\E}{\mathbb{E}}
\newcommand{\eps}{\varepsilon}
\newcommand{\dd}{\,\mathrm{d}}
\begin{document}
\section*{UnsolvedMath ID: 11000094}
\textbf{Source problem code:} catalogue entry
\subsection{Definitions and mathematical statement}
Retain the source's spaces \(\mathcal H_g\), \(\mathcal H_g^c\), their components \(\mathcal H_{g,\alpha}^c\), the target \(\mathfrak h_g\), and \(\mathcal J_g^c\), and its group notation \(T_g\). Their definitions are not included in the excerpt, so no identification with similarly named spaces or Lie algebras is presumed.

\emph{Problem 4.9.} Investigate the topology of \(\mathcal H_g\), \(\mathcal H_g^c\), and their components. Specifically, compute their homology and the cohomology at infinity of each \(\mathcal H_{g,\alpha}^c\). With a coefficient system \(R\) fixed, a usual ordinary-cohomology-at-infinity convention is
\[
 H^j_\infty(X;R)=\varinjlim_{K\subset X\ \mathrm{compact}}H^j(X\setminus K;R),
\]
with restriction maps as \(K\) increases. This is distinct from compactly supported cohomology. The source does not select coefficients or a convention for infinity, nor state the genus range or component indexing.

The source says the period map immerses each \(\mathcal H_{g,\alpha}^c\) in \(\mathfrak h_g\) as a closed subvariety; it concludes that each is a Stein manifold of complex dimension \(2g-1\), and hence has the homotopy type of a CW complex of that dimension. Those are retained as supplied structural claims, with definitions and the required smoothness/closed-map properties missing. Under the stated Stein-manifold hypothesis the usual CW consequence supplies an upper bound
\[
 \operatorname{hdim}(\mathcal H_{g,\alpha}^c)\leq2g-1,
\]
not by itself a lower bound or an assertion that every minimal CW model has dimension exactly \(2g-1\). The immersion/closed-subvariety wording must be checked before treating its Stein conclusion as established.

The adjacent numbered note says that the number of components of \(\mathcal H_g\) grows exponentially with genus while the minimum number of generators of \(T_g\) grows polynomially. It says the failure of an earlier genus-three argument in higher genus suggests that \(\mathcal J_g^c\) is generally not simply connected. The earlier argument, group definition, rates, and hypotheses are absent, and this suggestion is not converted into a proved non-simple-connectivity assertion.
\subsection{Short English statement}
Compute the homology of the source's \(\mathcal H_g\), \(\mathcal H_g^c\), and components, together with cohomology at infinity of \(\mathcal H_{g,\alpha}^c\). Its Stein and component-growth discussion is retained conditionally; the defining spaces and coefficient conventions are missing.
\subsection{Sources}
\begin{itemize}
\item Local UnsolvedMath registry, permanent ID 11000094, source code catalogue entry, Problem 4.9, PDF page 79. No individual proposer is supplied. Page identifies the item despite repeated chapter numbering.
\item Farb, \emph{Problems on Mapping Class Groups and Related Topics}. Supplied URL: \url{https://www.math.uchicago.edu/~farb/papers/mcgbook.pdf}; recorded access 2026-07-29; extraction pdf-font-structure. The edited volume contains 329 explicitly typeset headings across contributed chapters.
\item The statement supplies the period-map/Stein/CW discussion and a detached note numbered 10 about component growth and a genus-three argument, but no separate reference or earlier definitions.
\item Supplied status open and OPEN-TRIAGE. The research summary calls it a research program and says a settled answer cannot be verified. Current primary-literature status and redistribution terms for the publicly accessible author/publisher copy remain NEEDS\_REVIEW.
\end{itemize}
\end{document}
